% P11 paper module: R25 positive resolvent Jensen kernel and inverse-root strong criterion.

\subsection{Positive resolvent Jensen moments and the inverse-root fixed-vector criterion}
\label{subsec:o3x-resolvent-jensen}

Keep the odd-sector notation of Subsections~\ref{subsec:o3r-sqrt-offblock} and
\ref{subsec:o3w-resolvent-bridge} at fixed $0<R<S<T_0<U$:
\[
 A_R:=A_{T_0,U}^{R,-},
 \qquad
 A_S:=A_{T_0,U}^{S,-},
 \qquad
 W:=W_{R,S,-}^{[T_0]},
\]
\[
 \mathscr B_U:=A_SW-WA_R=(I-WW^*)A_SW.
\]
By Proposition~\ref{prop:o3w-coercivity},
\[
 A_R,A_S\ge a_0I,
 \qquad
 a_0=(1+\|H_{T_0}\|^2)^{-1}>0
\tag{O3X.1}\label{eq:o3x-a0}
\]
uniformly in $U$.

For $t\ge0$ put
\[
 R_R(t):=(A_R+tI)^{-1},
 \qquad
 R_S(t):=(A_S+tI)^{-1}.
\]

\begin{proposition}[Positive resolvent compression defect]
\label{prop:o3x-resolvent-gram}
Define
\[
 \mathscr E_U(t):=W^*R_S(t)W-R_R(t).
\]
Then
\[
 \boxed{
 \mathscr E_U(t)
 =R_R(t)\mathscr B_U^*R_S(t)\mathscr B_UR_R(t)
 \ge0.
 }
\tag{O3X.2}\label{eq:o3x-resolvent-gram}
\]
Equivalently,
\[
 \boxed{
 \mathscr E_U(t)
 =\bigl(R_S(t)^{1/2}\mathscr B_UR_R(t)\bigr)^*
  \bigl(R_S(t)^{1/2}\mathscr B_UR_R(t)\bigr).
 }
\tag{O3X.3}\label{eq:o3x-resolvent-square}
\]
\end{proposition}

\begin{proof}
The resolvent quasi-commutator is
\[
 R_S(t)W-WR_R(t)
 =-R_S(t)\mathscr B_UR_R(t).
\tag{O3X.4}\label{eq:o3x-resolvent-quasi}
\]
Taking adjoints and multiplying by $W$ gives
\[
 \mathscr E_U(t)
 =-R_R(t)\mathscr B_U^*R_S(t)W.
\]
Now insert~\eqref{eq:o3x-resolvent-quasi}.  Since
$\mathscr B_U^*W=0$, the $WR_R(t)$ term vanishes and one obtains
\eqref{eq:o3x-resolvent-gram}.  Formula~\eqref{eq:o3x-resolvent-square}
is immediate.
\end{proof}

The positive kernel in Proposition~\ref{prop:o3x-resolvent-gram} simultaneously
encodes the square-root Jensen defect from Section~\ref{sec:jensen}, a new logarithmic
defect, and the inverse-root defect of R24.

\begin{proposition}[Three Jensen moments of one positive kernel]
\label{prop:o3x-three-moments}
Let
\[
 \mathscr J_U:=A_R^{1/2}-W^*A_S^{1/2}W,
\]
so that $\mathscr J_U\ge0$ and
$\Theta=A_R^{-1/4}\mathscr J_UA_R^{-1/4}$ as in
\eqref{eq:jd-theta}.  Define
\[
 \mathscr L_U:=\log A_R-W^*(\log A_S)W,
\]
and
\[
 \mathscr I_U:=W^*A_S^{-1/2}W-A_R^{-1/2}.
\]
Then the following operator-norm Bochner integrals are exact:
\[
 \boxed{
 \mathscr J_U
 =\frac1\pi\int_0^\infty t^{1/2}\mathscr E_U(t)\,dt
 \ge0,
 }
\tag{O3X.5}\label{eq:o3x-sqrt-moment}
\]
\[
 \boxed{
 \mathscr L_U
 =\int_0^\infty\mathscr E_U(t)\,dt
 \ge0,
 }
\tag{O3X.6}\label{eq:o3x-log-moment}
\]
\[
 \boxed{
 \mathscr I_U
 =\frac1\pi\int_0^\infty t^{-1/2}\mathscr E_U(t)\,dt
 \ge0.
 }
\tag{O3X.7}\label{eq:o3x-inverse-moment}
\]
For every source vector $f$,
\[
 \boxed{
 \langle\mathscr L_Uf,f\rangle^2
 \le
 \pi^2\langle\mathscr J_Uf,f\rangle
       \langle\mathscr I_Uf,f\rangle.
 }
\tag{O3X.8}\label{eq:o3x-moment-interpolation}
\]
\end{proposition}

\begin{proof}
Equations~\eqref{eq:o3x-sqrt-moment} and
\eqref{eq:o3x-inverse-moment} follow from the standard positive functional-calculus
integrals for $A^{1/2}$ and $A^{-1/2}$ together with
Proposition~\ref{prop:o3x-resolvent-gram}.  For the logarithm use
\[
 \log A
 =\int_0^\infty
 \left(\frac1{1+t}I-(A+tI)^{-1}\right)dt;
\]
the scalar terms cancel after the isometric compression, giving
\eqref{eq:o3x-log-moment}.  Uniform coercivity handles the origin, while
$\|\mathscr E_U(t)\|=O(t^{-3})$ handles infinity, so all three integrals converge in
operator norm.

For fixed $f$ the measure
\[
 d\mu_{U,f}(t)
 :=\langle\mathscr E_U(t)f,f\rangle\,dt
\]
is positive.  Cauchy--Schwarz applied to
$1=t^{1/4}t^{-1/4}$ in $L^2(d\mu_{U,f})$ gives
\eqref{eq:o3x-moment-interpolation}.
\end{proof}

The inverse-root moment has a special advantage in the concrete P11 family: the
uniform lower spectral edge from R24 turns it into a cancellation-free fixed-vector
criterion.

\begin{theorem}[Positive fixed-vector criterion for inverse-root compatibility]
\label{thm:o3x-inverse-criterion}
Put
\[
 D_U^-:=A_S^{-1/2}W-WA_R^{-1/2}.
\]
Then, for every source vector $f$,
\[
 \boxed{
 \|D_U^-f\|^2
 \le a_0^{-1/2}\langle\mathscr I_Uf,f\rangle,
 }
\tag{O3X.9}\label{eq:o3x-inverse-vector-bound}
\]
and
\[
 \boxed{
 \mathscr I_U=W^*D_U^-.
 }
\tag{O3X.10}\label{eq:o3x-I-compression}
\]
Consequently, for every fixed $f$,
\[
 \boxed{
 D_U^-f\longrightarrow0
 \quad\Longleftrightarrow\quad
 \langle\mathscr I_Uf,f\rangle\longrightarrow0.
 }
\tag{O3X.11}\label{eq:o3x-fixed-equivalence}
\]
Moreover
\[
 \|D_U^-\|\le2a_0^{-1/2}
\tag{O3X.12}\label{eq:o3x-D-uniform}
\]
uniformly in $U$.  Hence it is enough to verify the scalar condition on any fixed
dense source core in order to obtain
\[
 D_U^-\xrightarrow[s]{}0
\]
on the whole source Hilbert space.
\end{theorem}

\begin{proof}
R24 gives
\[
 D_U^-
 =-\frac1\pi\int_0^\infty t^{-1/2}
 R_S(t)\mathscr B_UR_R(t)\,dt.
\]
By~\eqref{eq:o3x-resolvent-square} and
$R_S(t)^2\le(a_0+t)^{-1}R_S(t)$,
\[
 \|R_S(t)\mathscr B_UR_R(t)f\|^2
 \le(a_0+t)^{-1}
 \langle\mathscr E_U(t)f,f\rangle.
\]
Cauchy--Schwarz in the $t$-integral and
\[
 \int_0^\infty\frac{t^{-1/2}}{a_0+t}\,dt
 =\frac\pi{\sqrt{a_0}}
\]
prove~\eqref{eq:o3x-inverse-vector-bound}.  Equation
\eqref{eq:o3x-I-compression} follows directly from the definition of $D_U^-$.
The two implications in~\eqref{eq:o3x-fixed-equivalence} now follow from
\eqref{eq:o3x-inverse-vector-bound} and
\[
 |\langle\mathscr I_Uf,f\rangle|
 =|\langle D_U^-f,Wf\rangle|
 \le\|D_U^-f\|\,\|f\|.
\]
Finally~\eqref{eq:o3w-invroot-bound} gives
\eqref{eq:o3x-D-uniform}, and the standard dense-core approximation proves the strong
extension.
\end{proof}

\begin{proposition}[Scale-invariant logarithmic modulus defect]
\label{prop:o3x-log-defect}
The logarithmic defect is invariant under the common rescaling
$(A_R,A_S)\mapsto(cA_R,cA_S)$, $c>0$, and
\[
 \boxed{
 \mathscr L_U=0
 \quad\Longleftrightarrow\quad
 \mathscr B_U=0
 \quad\Longleftrightarrow\quad
 Q=W.
 }
\tag{O3X.13}\label{eq:o3x-log-zero}
\]
For the concrete P11 family there is $c>0$ such that, for all sufficiently large $U$,
\[
 \boxed{
 \|\mathscr L_U\|
 \ge c\,U^{-2m_h-2}.
 }
\tag{O3X.14}\label{eq:o3x-log-lower}
\]
\end{proposition}

\begin{proof}
Scale invariance follows from
$\log(cA)=(\log c)I+\log A$ and $W^*W=I$.  By
\eqref{eq:o3x-log-moment}, $\mathscr L_U=0$ iff the positive Gram kernel vanishes,
which is equivalent to $\mathscr B_U=0$.  For selfadjoint $A_S$, the latter means
$\Ran W$ reduces $A_S$, so functional calculus gives
$A_S^{1/2}W=WA_R^{1/2}$ and hence $Q=W$.  Conversely $Q=W$ gives that square-root
intertwining and, after squaring, $A_SW=WA_R$, hence $\mathscr B_U=0$.

For the quantitative statement put
\[
 M_U:=\max(\|A_R\|,\|A_S\|)=\|A_S\|.
\]
Choose a unit vector $f_U$ such that
$\|\mathscr B_Uf_U\|\ge\frac12\|\mathscr B_U\|$.  For $t\ge4M_U$,
\[
 \|R_R(t)-t^{-1}I\|\le M_Ut^{-2},
\]
so
\[
 \|\mathscr B_UR_R(t)f_U\|
 \ge\frac{\|\mathscr B_U\|}{4t}.
\]
Since $A_S\le M_UI$,
\[
 R_S(t)\ge(M_U+t)^{-1}I.
\]
Therefore
\[
 \langle\mathscr E_U(t)f_U,f_U\rangle
 \ge\frac{\|\mathscr B_U\|^2}{16t^2(M_U+t)}.
\]
Integration over $4M_U\le t\le8M_U$ gives a universal $c_*>0$ with
\[
 \|\mathscr L_U\|
 \ge c_*\frac{\|\mathscr B_U\|^2}{M_U^2}.
\tag{O3X.15}\label{eq:o3x-log-B-lower}
\]
Now use~\eqref{eq:o3l-B-lower} and
Lemma~\ref{lem:o3l-relative-upper} to obtain
\eqref{eq:o3x-log-lower}.
\end{proof}

\begin{remark}[R25 firewall: positive resolvent moments remain modulus data]
\label{rem:o3x-firewall}
The new positivity removes cancellation in the spectral parameter but does not cross the
polar-gauge firewall.  In the R14 perfect-modulus model,
$Q=W$ and therefore
\[
 \mathscr B_U=\mathscr E_U(t)=\mathscr L_U=\mathscr I_U=0
\]
while the actual polar gauge remains nontrivial.  Conversely the R19 and R20 combined
models can carry nonzero modulus/Jensen defects while their relative polar defect
vanishes.  Thus no estimate in this subsection alone implies a bound on
$U_SW-WU_R$ or on the R22 angle defect $\mathscr G_U$.
\end{remark}

\begin{openproblem}[R25-F: positive inverse-root defect measure]
\label{open:o3x-inverse-measure}
For each fixed smooth odd source vector $f$, determine whether
\[
 \boxed{
 \langle\mathscr I_Uf,f\rangle
 =\frac1\pi\int_0^\infty t^{-1/2}
 \|R_S(t)^{1/2}\mathscr B_UR_R(t)f\|^2\,dt
 \longrightarrow0
 }
\]
or construct a fixed $f$ and a subsequence with positive limsup.  By
Theorem~\ref{thm:o3x-inverse-criterion}, this is exactly the fixed-vector version of the
R24 inverse-root compatibility problem and contains no cancellation in $t$.

A resolution of this problem is still a modulus/inverse-root result.  Separate relative
polar information remains necessary for R22-F and for strong terminal transport.
\end{openproblem}
